\documentclass[11pt]{article}
\usepackage[margin=1in]{geometry}
\usepackage{amsmath,amssymb,amsthm,amsfonts}
\usepackage{mathtools}
\usepackage{bm}
\usepackage{hyperref}
\usepackage{enumitem}
\usepackage{graphicx}
\usepackage{bbm}
\usepackage{booktabs}
\usepackage{microtype}

\title{Neuromorphic AGI, Multiplicity Theory, and the Universal Multiplicity Equation:\\
From Descriptive PDE to Operator Generator and Executable Simulator}
\author{Draft Synthesis based on April 2026 Discussion}
\date{\today}

\begin{document}
\maketitle

\begin{abstract}
This report synthesizes a sequence of developments around Multiplicity Theory, Neuromorphic Artificial General Intelligence (AGI), and the Universal Multiplicity Equation (UME). The original neuromorphic AGI preprint presents a time-dependent multiplicity equation as a unifying dynamical law for cognitive state evolution, integrating prime-based encoding, recursive feedback, tensor interactions, diffusion, and stochasticity. Here we upgrade that formulation in three directions: (1) we re-cast the UME from a descriptive partial differential equation into a generator of a semigroup on a multiplicity Hilbert space; (2) we formalize prime indices as labels for irreducible operator families acting on this space; and (3) we provide a concrete PyTorch-style architecture (the \emph{MultiplicityCell}) together with an experiment design to test a prime-factor transfer law across composite tasks. The report concludes with recommendations for next-step experiments and validation pathways.
\end{abstract}

\newpage
\tableofcontents


\section{Executive Summary}

\subsection{Starting point: neuromorphic UME}

The neuromorphic AGI preprint introduces a Universal Multiplicity Equation
\begin{equation}
\frac{\partial \psi_k(t)}{\partial t}
= \alpha_k(t)\psi_k
+ \beta_k(t)\int_0^t I_k(\tau)\,d\tau
+ \gamma_k(t)\sum_{j,l} T_{kjl}\psi_j\psi_l
+ \lambda_k(t)\nabla^2\psi_k
+ \eta_k(t)\psi_k^n
+ \xi_k(t)
\label{eq:UME-coordinate}
\end{equation}
with explicit dimensional analysis and characteristic scales $T_s$ and $L_s$ ensuring that each term matches the dimensions of the left-hand side.\footnote{All references to the neuromorphic AGI preprint and its equations are drawn from the attached document.} This equation is presented as a framework for cognitive state evolution in neuromorphic AGI, combining:
\begin{itemize}[nosep]
\item Growth/decay ($\alpha_k$),
\item Memory integrals ($\beta_k$, $I_k$),
\item Tensor-mediated coupling ($T_{kjl}$),
\item Diffusion ($\lambda_k\nabla^2$),
\item Nonlinear self-interaction ($\eta_k\psi_k^n$),
\item Stochasticity ($\xi_k$).
\end{itemize}

The preprint further incorporates quantum-inspired structure via prime-encoded superpositions, tensorial entanglement, and geometric feedback terms using Berry curvature and the Fubini--Study metric.

\subsection{Conceptual upgrade}

The key upgrade developed in the discussion is to reinterpret \eqref{eq:UME-coordinate} not merely as a PDE written in components, but as the coordinate representation of an operator equation
\begin{equation}
\boxed{
\frac{d\Psi}{dt} = \mathcal{L}_{\mathrm{mult}}[\Psi]
}
\end{equation}
where $\Psi$ is a state in a (possibly infinite-dimensional) multiplicity Hilbert space $\mathcal{H}$, and $\mathcal{L}_{\mathrm{mult}}$ is a generator of a strongly continuous semigroup (or group, when unitary) incorporating distinct operator contributions.

The operator decomposition is schematically
\begin{equation}
\mathcal{L}_{\mathrm{mult}}
=
\underbrace{\Xi(t)}_{\text{prime evolution}}
+
\underbrace{\Lambda_m(t)}_{\text{multiplicity scalar}}
+
\underbrace{\mathcal{T}}_{\text{tensor coupling}}
+
\underbrace{\mathcal{D}}_{\text{diffusion}}
+
\underbrace{\mathcal{N}}_{\text{nonlinear feedback}}
+
\underbrace{\mathcal{S}}_{\text{stochastic term}}.
\end{equation}

Prime indices are elevated from static labels to indices of irreducible operator families $\{\hat{O}_p\}$; composite cognition is modeled as operator words $W = \prod_p \hat{O}_p^{k_p}$ acting on $\Psi$. Ethics is upgraded from a linear update rule to a geometric constraint manifold $\mathcal{E}\subset\mathcal{H}$ defined by a fairness functional $F[\Psi]\ge \varepsilon$.

\subsection{Executable upgrade: MultiplicityCell and transfer experiment}

We then construct a concrete \emph{MultiplicityCell} in PyTorch that discretizes the dynamics
\[
\Psi_{t+1} = \Psi_t + \Delta t\,\mathcal{L}_{\mathrm{mult}}(\Psi_t)
\]
over a discrete multiplicity axis representing prime-encoded basis states. The cell includes:
\begin{enumerate}[nosep]
\item A diagonal prime-evolution operator,
\item A nonlinear multiplicity-scalar operator,
\item A tensor-like mixing operator across prime indices,
\item A discrete Laplacian as diffusion operator,
\item Nonlinear self-interaction,
\item Optional ethical projection and geometric regularization via Fubini--Study curvature penalties.
\end{enumerate}

On top of this cell, we design a multi-task experiment to test a \emph{prime-factor transfer law}. Tasks are assigned to subsets of prime indices; composite tasks use unions of these subsets. The law predicts that transfer performance from base tasks to a composite task scales with the normalized overlap of their activated prime sets. This is contrasted with random partition baselines and non-prime-structured RNNs.

\section{The Universal Multiplicity Equation: Coordinate Form}

The neuromorphic AGI preprint presents the UME in the following form (here we adopt their notation, suppressing explicit scaling factors for readability):
\begin{equation}
\frac{\partial \psi_k(t)}{\partial t}
=
\alpha_k(t)\psi_k
+
\beta_k(t)\int_0^t I_k(\tau)\,d\tau
+
\gamma_k(t)\sum_{j,l} T_{kjl}\psi_j\psi_l
+
\lambda_k(t)\nabla^2\psi_k
+
\eta_k(t)\psi_k^n
+
\xi_k(t).
\end{equation}
Key points:
\begin{itemize}[nosep]
\item Each index $k$ indexes a cognitive mode, which in the multiplicity picture will be associated with a prime-labeled module.
\item The coefficients $\alpha_k, \beta_k, \gamma_k, \lambda_k, \eta_k$ may themselves be time-dependent and subject to learning.
\item The Laplacian $\nabla^2$ acts on spatial or graph coordinates associated with neuromorphic substrates.
\item The integral term $\int_0^t I_k(\tau)\,d\tau$ captures history-dependent effects, aligning with memory and synaptic plasticity.
\end{itemize}

In addition, the preprint embeds quantum-inspired structure:
\begin{itemize}[nosep]
\item Superposition of prime-encoded states:
  \[
  \psi(t) = \sum_{k=1}^N c_k \lvert p_k \rangle
  \]
  with $c_k$ complex amplitudes and $\lvert p_k\rangle$ prime-labelled basis states.
\item Entanglement via tensor terms $T_{kjl}$.
\item Geometric feedback via Berry curvature $\Omega_B$ and the Fubini--Study metric $\Omega_{\mathrm{FS}}$.
\end{itemize}

Ethical adaptation appears in discrete-time form as
\begin{equation}
M(t+1) = \alpha M(t) + \beta R(t),
\end{equation}
and prime-weighted ethical weights are summarized by an expression of the form
\begin{equation}
P_{\mathrm{ethical}} = \sum_{k=1}^N \frac{1}{p_k},
\end{equation}
where $p_k$ are primes associated with ethical weightings.

\section{Operator Algebra Upgrade}

\subsection{Multiplicity Hilbert space}

We introduce a multiplicity Hilbert space $\mathcal{H}$ which, in a simple realization, could be
\[
\mathcal{H} \cong \ell^2(\mathbb{P}) \otimes \mathcal{K}
\]
where:
\begin{itemize}[nosep]
\item $\ell^2(\mathbb{P})$ is the square-summable space over primes, with canonical basis $\{\lvert p_k\rangle\}$,
\item $\mathcal{K}$ is a feature space (e.g., $L^2(\Omega)$ over spatial domain $\Omega$, or $\mathbb{C}^d$ for finite feature vectors).
\end{itemize}

A general state can then be written as
\[
\Psi = \sum_{k} \psi_k \otimes \lvert p_k\rangle,
\]
with $\psi_k \in \mathcal{K}$ corresponding to the component associated with prime $p_k$.

\subsection{Operator decomposition of $\mathcal{L}_{\mathrm{mult}}$}

We define $\mathcal{L}_{\mathrm{mult}}$ as an operator (possibly time-dependent) on $\mathcal{H}$, decomposed as
\begin{equation}
\mathcal{L}_{\mathrm{mult}} =
\Xi(t) + \Lambda_m(t) + \mathcal{T} + \mathcal{D} + \mathcal{N} + \mathcal{S}.
\end{equation}

\paragraph{Prime evolution operator $\Xi(t)$}

We let
\[
\Xi(t) = \sum_{k} \alpha_k(t)\, \mathbb{I} \otimes \lvert p_k\rangle\langle p_k\rvert,
\]
which acts diagonally on the prime index, matching the growth/decay term $\alpha_k(t)\psi_k$.

\paragraph{Multiplicity-scalar operator $\Lambda_m(t)$}

This is a nonlinear or linear operator acting within each prime block, e.g.,
\[
\Lambda_m(t)[\Psi] = \sum_k \Lambda^{(k)}_m(t)\psi_k \otimes \lvert p_k\rangle,
\]
where $\Lambda_m^{(k)}(t)$ may be realized as a small neural network or MLP on the feature space $\mathcal{K}$.

\paragraph{Tensor coupling operator $\mathcal{T}$}

We interpret the $T_{kjl}$ tensor as defining an operator-valued quadratic form:
\begin{equation}
\mathcal{T}[\Psi] = \sum_{k,j,l} T_{kjl}\,(\psi_j,\psi_l)\otimes \lvert p_k\rangle,
\end{equation}
where $(\psi_j,\psi_l)$ is a suitable bilinear combination (e.g., pointwise product, feature-wise product, or more general bilinear form). In a low-rank factorization,
\[
T_{kjl} = \sum_{r=1}^R U_{kr}V_{jr}W_{lr},
\]
leading to efficient tensor network implementations.

\paragraph{Diffusion operator $\mathcal{D}$}

We let
\[
\mathcal{D} = \sum_k \lambda_k(t) \Delta \otimes \lvert p_k\rangle\langle p_k\rvert,
\]
where $\Delta$ is a Laplacian on $\mathcal{K}$ (spatial, graph, or otherwise). On a discrete prime axis, we may also consider a Laplacian in the prime index itself.

\paragraph{Nonlinear operator $\mathcal{N}$}

The nonlinear self-interaction term $\eta_k(t)\psi_k^n$ is captured as
\[
\mathcal{N}[\Psi] = \sum_k \eta_k(t)\, f_n(\psi_k)\otimes \lvert p_k\rangle,
\]
where $f_n$ denotes a power or general nonlinear map.

\paragraph{Stochastic term $\mathcal{S}$}

The stochastic term $\xi_k(t)$ is modeled as noise:
\[
\mathcal{S}[\Psi,t] = \Xi_{\mathrm{noise}}(t),
\]
with $\Xi_{\mathrm{noise}}(t)$ a stochastic process in $\mathcal{H}$ (e.g., Gaussian in each coordinate).

\subsection{Prime operators as symmetry generators}

Instead of treating primes as mere labels, we introduce a family of operators $\{\hat{O}_p\}$ on $\mathcal{H}$, each associated with a prime $p$. Composite cognition is then an operator word
\[
W = \prod_p \hat{O}_p^{k_p},
\]
acting on $\Psi$. A task or cognitive transformation can be represented as such an operator word, or as a word in a noncommutative algebra generated by the $\hat{O}_p$:
\[
[\hat{O}_p,\hat{O}_q]\neq 0
\]
whenever cognitive primitives associated with $p$ and $q$ interact non-trivially.

\subsection{Ethics as constraint manifold}

We upgrade the discrete ethical update
\[
M(t+1) = \alpha M(t) + \beta R(t)
\]
to a geometric constraint:
\begin{equation}
\Psi(t)\in\mathcal{E}\subset\mathcal{H},
\end{equation}
with
\[
\mathcal{E} = \{\Psi\in\mathcal{H} \mid F[\Psi] \geq \varepsilon\},
\]
where $F$ is a fairness or lawfulness functional. The dynamics are then restricted to $\mathcal{E}$ either by:
\begin{itemize}[nosep]
\item projected evolution, $\Psi_{t+1} = \Pi_{\mathcal{E}}\big(\Psi_t + \Delta t\,\mathcal{L}_{\mathrm{mult}}(\Psi_t)\big)$, or
\item penalization in a variational functional, e.g.,
\[
\mathcal{F}[\Psi] = E[\Psi] + R[\Psi] + C[\Psi] + \lambda_{\mathcal{E}}\max(0,\varepsilon - F[\Psi]).
\]
\end{itemize}

\subsection{Quantum geometry: Fubini--Study and curvature}

The preprint references the Fubini--Study metric and Berry curvature as mechanisms for geometric feedback. On the projective Hilbert space $\mathbb{P}(\mathcal{H})$, the Fubini--Study distance between normalized states $\Psi_1,\Psi_2$ is
\[
d_{\mathrm{FS}}(\Psi_1,\Psi_2)
=
\arccos\big(|\langle \Psi_1,\Psi_2\rangle|\big).
\]
A curvature penalty can be introduced over trajectories $\Psi(t)$ by integrating a curvature density $\kappa(\Psi(t))$, or, in a discrete simulation, by penalizing large Fubini--Study distances between successive states:
\[
C[\{\Psi_t\}] = \sum_t d_{\mathrm{FS}}(\Psi_t,\Psi_{t+1}).
\]
This yields a concrete geometric regularization term compatible with the geometric feedback concept.

\section{Variational Learning and Prime-Factor Transfer Law}

\subsection{Variational principle}

We propose to treat learning as a constrained variational problem:
\begin{equation}
\Psi^\ast = \arg\min_{\Psi \in \mathcal{E}} \mathcal{F}[\Psi],
\end{equation}
where
\begin{equation}
\mathcal{F} = E[\Psi] + R[\Psi] + C[\Psi],
\end{equation}
with:
\begin{itemize}[nosep]
\item $E[\Psi]$: stability or energy functional (e.g., norms or Lyapunov-like function),
\item $R[\Psi]$: task loss (e.g., prediction error),
\item $C[\Psi]$: curvature/geometry penalty (e.g., Fubini--Study regularization).
\end{itemize}
Ethical constraints enter via $\mathcal{E}$ or an additional penalty term as above.

\subsection{Prime-factor transfer law}

Let each task $T_i$ be associated with a subset $P_i$ of prime indices (or their operator families). We define a normalized overlap
\[
\mathrm{Overlap}(i,j) = \frac{|P_i\cap P_j|}{|P_i\cup P_j|}.
\]
The prime-factor transfer law posits that transfer performance from $T_i$ to $T_j$ satisfies
\begin{equation}
\mathrm{Transfer}(T_i\to T_j)
\approx
a\,\mathrm{Overlap}(i,j) + b,
\end{equation}
for constants $a>0$, $b$, up to noise. Transfer can be measured via zero-shot performance, few-shot sample efficiency, or reduction in loss after a fixed adaptation budget.

\section{Discrete Simulator: MultiplicityCell}

We discretize the UME dynamics over time with step size $\Delta t$, on a finite-dimensional approximation of $\mathcal{H}$ given by tensors of shape
\[
\text{batch} \times \text{feature\_dim} \times \text{multiplicity\_dim}.
\]
Here, \texttt{multiplicity\_dim} indexes prime basis states $\lvert p_k\rangle$.

\subsection{Cell definition}

Below is a PyTorch-style definition of one time-step of the discretized dynamics:
\begin{verbatim}
import torch
import torch.nn as nn
import torch.nn.functional as F

class MultiplicityCell(nn.Module):
    """
    One step of UME dynamics:
        dΨ/dt = L_mult[Ψ]
    Discretized:
        Ψ_{t+1} = Ψ_t + Δt * L_mult(Ψ_t)
    Ψ has shape (batch, feature_dim, multiplicity_dim).
    """

    def __init__(self, feature_dim, multiplicity_dim,
                 hidden_dim, delta_t=0.1, nonlinear_order=3):
        super().__init__()
        self.feature_dim = feature_dim
        self.multiplicity_dim = multiplicity_dim
        self.hidden_dim = hidden_dim
        self.delta_t = delta_t
        self.nonlinear_order = nonlinear_order

        # 1. Prime evolution operator α_k (diagonal in multiplicity basis)
        self.alpha = nn.Parameter(torch.randn(multiplicity_dim) * 0.1)

        # 2. Multiplicity-scalar operator (nonlinear per-prime feature transform)
        self.scalar_net = nn.Sequential(
            nn.Linear(feature_dim, hidden_dim),
            nn.ReLU(),
            nn.Linear(hidden_dim, feature_dim)
        )

        # 3. Tensor-like coupling operator (mixing across multiplicities)
        self.mixing_matrix = nn.Parameter(
            torch.randn(multiplicity_dim, multiplicity_dim) * 0.1
        )
        self.coupling_proj = nn.Linear(feature_dim, feature_dim)

        # 4. Diffusion operator λ_k ∇² over multiplicity axis
        self.lambda_diff = nn.Parameter(torch.randn(multiplicity_dim) * 0.1)
        self.register_buffer("laplacian",
                             self._compute_laplacian(multiplicity_dim))

        # 5. Nonlinear self-interaction η_k ψ_k^n
        self.nonlinear_eta = nn.Parameter(torch.randn(multiplicity_dim) * 0.01)

        # 6. Optional ethical projection (constraint manifold)
        self.ethical_projection = None

    def _compute_laplacian(self, n):
        """Simple 1D graph Laplacian on multiplicity indices."""
        L = torch.diag(torch.full((n,), -2.0))
        for i in range(n - 1):
            L[i, i + 1] = 1.0
            L[i + 1, i] = 1.0
        return L

    def forward(self, psi, external_input=None):
        """
        psi: (batch, feature_dim, multiplicity_dim)
        external_input: optional (batch, feature_dim, multiplicity_dim)
        """
        batch_size = psi.size(0)

        # Term 1: prime evolution (α_k ψ_k)
        term1 = self.alpha.view(1, 1, -1) * psi

        # Term 2: multiplicity-scalar nonlinearity per prime
        psi_flat = psi.permute(0, 2, 1).reshape(-1, self.feature_dim)
        scalar_out = self.scalar_net(psi_flat)
        term2 = scalar_out.reshape(batch_size, self.multiplicity_dim,
                                   self.feature_dim).permute(0, 2, 1)

        # Term 3: tensor-like coupling via mixing across multiplicities
        # mixing_matrix: (k, j), psi: (b, f, j) -> (b, f, k)
        mixed = torch.einsum('kj, bfj -> bfk', self.mixing_matrix, psi)
        term3 = self.coupling_proj(mixed.permute(0, 2, 1)).permute(0, 2, 1)

        # Term 4: diffusion over multiplicity axis
        lap_psi = torch.einsum('kj, bfj -> bfk', self.laplacian, psi)
        term4 = self.lambda_diff.view(1, 1, -1) * lap_psi

        # Term 5: nonlinear self-interaction
        term5 = self.nonlinear_eta.view(1, 1, -1) * (psi ** self.nonlinear_order)

        # External input / stochastic term ξ(t)
        if external_input is None:
            noise = torch.randn_like(psi) * 0.01
        else:
            noise = external_input

        L_mult = term1 + term2 + term3 + term4 + term5 + noise

        psi_next = psi + self.delta_t * L_mult

        # Optional ethical projection (constraint manifold)
        if self.ethical_projection is not None:
            psi_next = self.ethical_projection(psi_next)

        return psi_next
\end{verbatim}

\subsection{Geometric regularization}

We implement a discrete Fubini--Study proxy and curvature penalty as:
\begin{verbatim}
def compute_fs_distance(psi1, psi2):
    b = psi1.size(0)
    v1 = psi1.view(b, -1)
    v2 = psi2.view(b, -1)
    v1 = F.normalize(v1, dim=1)
    v2 = F.normalize(v2, dim=1)
    overlap = torch.abs(torch.sum(v1 * v2, dim=1))
    overlap = torch.clamp(overlap, -0.999, 0.999)
    return torch.acos(overlap)

def curvature_penalty(psi_seq):
    total = 0.0
    for t in range(len(psi_seq) - 1):
        dist = compute_fs_distance(psi_seq[t], psi_seq[t + 1])
        total += dist.mean()
    return total
\end{verbatim}

This implements the geometric feedback notion referenced in the UME framework in a form directly usable in gradient-based learning.

\section{Prime-Structured Multi-Task Transfer Experiment}

\subsection{Task family}

We consider a family of tasks, each defined as an operator composition acting on input vectors in a feature space. For a toy example, we may define primitives:

\begin{center}
\begin{tabular}{lll}
\toprule
Prime index & Primitive & Description \\
\midrule
0 & identity      & $x \mapsto x$ \\
1 & negate        & $x \mapsto -x$ (subset of components) \\
2 & double        & $x \mapsto 2x$ (subset) \\
3 & square        & $x \mapsto x^{\odot 2}$ \\
4 & shift\_left   & circular shift of features \\
5 & threshold     & elementwise step or ReLU-like map \\
6 & smooth        & local averaging across neighboring features \\
7 & noise\_gate   & dropout-like gating \\
\bottomrule
\end{tabular}
\end{center}

Tasks are then defined as compositions $T_A = \hat{O}_1\circ\hat{O}_2$, $T_B = \hat{O}_3$, $T_C = \hat{O}_4\circ\hat{O}_6$, etc., with associated prime support sets $P_A, P_B, P_C\subset\{0,\dots,7\}$. Composite tasks such as $T_D = \hat{O}_1\circ\hat{O}_3\circ\hat{O}_4$ then have $P_D = P_A\cup P_B\cup P_C$ where appropriate.

\subsection{Training protocol}

We propose:
\begin{enumerate}[nosep]
\item Pre-train a model built from the MultiplicityCell on base tasks $A,B,C$ under prime masks that restrict learning to their corresponding prime subsets.
\item Freeze most of the model (except perhaps a small readout head).
\item Evaluate zero-shot performance on composite task $D$ and a disjoint-overlap control task $E$.
\item Optionally fine-tune with a small number of gradient steps (few-shot regime) and measure sample efficiency and interference with base tasks.
\end{enumerate}

\subsection{Prime-mask vs random-mask controls}

To test whether prime-structured modularity is doing real work, we introduce:
\begin{itemize}[nosep]
\item A random-partition baseline, where task-specific masks are sampled at random but with matching cardinalities and overlap statistics.
\item A non-masked RNN baseline (same parameter budget, no explicit multiplicity structure).
\end{itemize}

We then compare transfer metrics across conditions.

\subsection{Metrics and expected outcomes}

Key metrics:
\begin{itemize}[nosep]
\item Zero-shot accuracy/error on composite tasks.
\item Few-shot learning curves (loss vs. steps).
\item Catastrophic interference on base tasks after fine-tuning on composite tasks.
\item Internal diagnostics: per-prime activation energy, coupling usage, Fubini--Study distances between task trajectories.
\end{itemize}

The prime-factor transfer law predicts that
\begin{equation}
\mathrm{Transfer}(T_i\to T_j)
\propto \frac{|P_i\cap P_j|}{|P_i\cup P_j|}
\end{equation}
and that prime-structured masks outperform random partitions and unstructured RNNs, controlling for resource budget.

\section{Conclusions and Recommended Next Steps}

This report has:
\begin{enumerate}[nosep]
\item Recast the neuromorphic UME from a descriptive PDE into an operator-generator framework on multiplicity space.
\item Clarified the role of primes as indices of operator families and compositional modes.
\item Upgraded the ethical component from a linear update rule to a geometric constraint manifold.
\item Instantiated a concrete MultiplicityCell architecture that discretizes the UME in a differentiable, trainable form.
\item Defined an experiment to test a prime-factor transfer law across tasks with controlled prime-support overlap.
\end{enumerate}

Immediate next steps that would most increase the scientific weight of the framework are:
\begin{itemize}[nosep]
\item Implement and run the prime-factor transfer experiment with overlap-matched random controls and non-masked RNN baselines.
\item Quantitatively fit the transfer law and report confidence intervals on the overlap--performance correlation.
\item Extend the simulator to sequence tasks and, eventually, spiking/neuromorphic substrates while preserving the multiplicity semantics.
\end{itemize}

These steps would turn the current conceptual unification into an empirically grounded, falsifiable ``standard model'' of multiplicity-governed cognition.

\appendix
\section*{Mathematical Appendix}

\section{Abstract Operator Setting and Basic Bounds}

\subsection{State space and notation}

Let \(\mathbb{P}_N = \{p_1,\dots,p_N\}\) be the first \(N\) primes and let
\[
\mathcal{H} \coloneqq \mathcal{K} \otimes \ell^2(\mathbb{P}_N)
\]
be the finite-dimensional multiplicity Hilbert space, where \(\mathcal{K} \cong \mathbb{C}^d\) is a feature space and \(\ell^2(\mathbb{P}_N)\) is spanned by \(\{\lvert p_k\rangle\}_{k=1}^N\). A state \(\Psi\in\mathcal{H}\) can be written as
\[
\Psi = \sum_{k=1}^N \psi_k \otimes \lvert p_k\rangle,
\quad \psi_k \in \mathcal{K}.
\]

We equip \(\mathcal{H}\) with the usual Hilbert norm
\[
\|\Psi\|_{\mathcal{H}}^2
=
\sum_{k=1}^N \|\psi_k\|_{\mathcal{K}}^2.
\]

In this appendix, all operators act on the finite-dimensional \(\mathcal{H}\); all operator norms \(\|\cdot\|\) are taken in the induced operator norm on \(\mathcal{H}\).

\subsection{Operator decomposition of \(\mathcal{L}_{\mathrm{mult}}\)}

We consider an operator of the form
\begin{equation}
\mathcal{L}_{\mathrm{mult}}
=
\Xi + \Lambda_m + \mathcal{T} + \mathcal{D} + \mathcal{N} + \mathcal{S},
\label{eq:Lmult-decomp}
\end{equation}
where each term corresponds to one of the dynamical contributions appearing in the coordinate UME, but now lifted to \(\mathcal{H}\). For the purposes of this appendix we restrict to the deterministic part and treat \(\mathcal{S}\) as either zero or a bounded perturbation; stochastic analysis can be added on top of these bounds if desired.

We define each operator and derive explicit norm bounds.

\subsection{Prime evolution operator \(\Xi\)}

Define
\begin{equation}
\Xi
\coloneqq
\sum_{k=1}^N \alpha_k \left(\mathbb{I}_{\mathcal{K}} \otimes \lvert p_k\rangle\langle p_k\rvert\right),
\end{equation}
where \(\alpha_k \in \mathbb{C}\) are scalar coefficients. Then \(\Xi\) acts diagonally:
\[
(\Xi \Psi)_k = \alpha_k \psi_k.
\]

\begin{lemma}
\label{lem:Xi-norm}
The operator norm of \(\Xi\) satisfies
\[
\|\Xi\| = \max_{1\le k\le N} |\alpha_k|.
\]
\end{lemma}

\begin{proof}
For any \(\Psi = \sum_k \psi_k \otimes \lvert p_k\rangle\),
\[
\|\Xi \Psi\|_{\mathcal{H}}^2
=
\sum_{k=1}^N \|\alpha_k \psi_k\|_{\mathcal{K}}^2
=
\sum_{k=1}^N |\alpha_k|^2 \|\psi_k\|_{\mathcal{K}}^2
\le
\Bigl(\max_k |\alpha_k|^2\Bigr)\sum_k \|\psi_k\|_{\mathcal{K}}^2
=
\bigl(\max_k|\alpha_k|\bigr)^2 \|\Psi\|_{\mathcal{H}}^2.
\]
Thus \(\|\Xi\|\le \max_k|\alpha_k|\). Equality is achieved by taking \(\Psi\) supported entirely in the index \(k^\ast\) with \(|\alpha_{k^\ast}|=\max_k|\alpha_k|\), so \(\|\Xi\| = \max_k|\alpha_k|\).
\end{proof}

\subsection{Multiplicity-scalar operator \(\Lambda_m\)}

We model \(\Lambda_m\) as a block-diagonal linear map on \(\mathcal{H}\):
\begin{equation}
\Lambda_m \Psi
\coloneqq
\sum_{k=1}^N (B_k \psi_k)\otimes \lvert p_k\rangle,
\end{equation}
where each \(B_k:\mathcal{K}\to \mathcal{K}\) is a bounded linear operator (for example, the linearization of a small MLP or a learned per-prime linear map).

\begin{lemma}
\label{lem:Lm-norm}
If \(\|B_k\| \le b_k\) for all \(k\), then
\[
\|\Lambda_m\| \le \max_{1\le k\le N} b_k.
\]
In particular, if \(\|B_k\|\le b\) uniformly, then \(\|\Lambda_m\|\le b\).
\end{lemma}

\begin{proof}
For any \(\Psi\in\mathcal{H}\),
\[
\|\Lambda_m \Psi\|_{\mathcal{H}}^2
=
\sum_{k=1}^N \|B_k \psi_k\|_{\mathcal{K}}^2
\le
\sum_{k=1}^N b_k^2 \|\psi_k\|_{\mathcal{K}}^2
\le
\bigl(\max_k b_k^2\bigr)\sum_k \|\psi_k\|_{\mathcal{K}}^2
=
\bigl(\max_k b_k\bigr)^2 \|\Psi\|_{\mathcal{H}}^2.
\]
The conclusion follows.
\end{proof}

\subsection{Diffusion operator \(\mathcal{D}\)}

Let \(L \in \mathbb{C}^{N\times N}\) be a graph Laplacian on the prime index set. In the simplest case used in the cell, \(L\) is the 1D chain Laplacian with entries
\[
L_{ij} =
\begin{cases}
-2, & i=j \in \{2,\dots,N-1\},\\
-1, & i=j\in\{1,N\},\\
1,  & |i-j|=1,\\
0,  & \text{otherwise.}
\end{cases}
\]
Define
\begin{equation}
(\mathcal{D}\Psi)_k
\coloneqq
\lambda_k \sum_{j=1}^N L_{kj}\psi_j,
\end{equation}
where \(\lambda_k\in \mathbb{C}\). In matrix form this is
\[
\mathcal{D} = (\Lambda_\lambda \otimes \mathbb{I}_{\mathcal{K}})\,(\mathbb{I}_{\mathcal{K}}\otimes L),
\]
with \(\Lambda_\lambda = \mathrm{diag}(\lambda_1,\dots,\lambda_N)\).

\begin{lemma}
\label{lem:D-norm}
The operator norm of \(\mathcal{D}\) is bounded by
\[
\|\mathcal{D}\|
\le
\bigl(\max_{k} |\lambda_k|\bigr) \|L\|_{2},
\]
where \(\|L\|_2\) is the spectral norm of the Laplacian on \(\mathbb{C}^N\).
\end{lemma}

\begin{proof}
We view \(\mathcal{H}\cong \mathcal{K}^N\) as a direct sum and note that
\[
\mathcal{D}
=
(\Lambda_\lambda \otimes \mathbb{I}_{\mathcal{K}})\,(\mathbb{I}_{\mathcal{K}}\otimes L),
\]
so
\[
\|\mathcal{D}\|
\le
\|\Lambda_\lambda\otimes \mathbb{I}_{\mathcal{K}}\|\,
\|\mathbb{I}_{\mathcal{K}}\otimes L\|
=
\|\Lambda_\lambda\|\,\|L\|
=
\bigl(\max_k|\lambda_k|\bigr)\,\|L\|_2,
\]
using standard properties of Kronecker products and the fact that \(\|\mathbb{I}_{\mathcal{K}}\|=1\).
\end{proof}

\subsection{Tensor coupling operator \(\mathcal{T}\)}

We model the tensor coupling in a bounded linear approximation. Let \(C\in\mathbb{C}^{N\times N}\) be a mixing matrix and \(W:\mathcal{K}\to\mathcal{K}\) a bounded linear operator. Define
\begin{equation}
(\mathcal{T}\Psi)_k \coloneqq W\left(\sum_{j=1}^N C_{kj}\psi_j\right).
\end{equation}
This matches the form used in the cell, where \(C\) is the learnable mixing matrix and \(W\) corresponds to the coupling projection.

\begin{lemma}
\label{lem:T-norm}
If \(\|W\|\le w\) and \(\|C\|_2\le c\) (spectral norm), then
\[
\|\mathcal{T}\|\le wc.
\]
\end{lemma}

\begin{proof}
Identify \(\mathcal{H}\cong \mathcal{K}^N\). The map \(\Psi\mapsto \sum_j C_{kj}\psi_j\) is \(\mathbb{I}_{\mathcal{K}}\otimes C\) on \(\mathcal{H}\), and the subsequent application of \(W\) at each index is \(W\otimes \mathbb{I}_{\ell^2(\mathbb{P}_N)}\). Thus
\[
\mathcal{T} = (W\otimes \mathbb{I})\,(\mathbb{I}\otimes C),
\]
and
\[
\|\mathcal{T}\|
\le
\|W\otimes \mathbb{I}\|\,\|\mathbb{I}\otimes C\|
=
\|W\|\,\|C\|
\le wc.
\]
\end{proof}

\subsection{Nonlinear operator \(\mathcal{N}\)}

We consider a polynomial-type nonlinearity per prime:
\begin{equation}
(\mathcal{N}\Psi)_k
\coloneqq
\eta_k f_n(\psi_k)
\end{equation}
where \(f_n:\mathcal{K}\to\mathcal{K}\) is defined componentwise by \(f_n(x)=x^{\odot n}\) in a chosen basis and \(\eta_k\in\mathbb{C}\). In finite dimension, \(\mathcal{N}\) is locally Lipschitz.

\begin{lemma}[Local Lipschitz bound]
\label{lem:N-local}
Fix \(R>0\). On the ball \(\{\Psi\in\mathcal{H} : \|\Psi\|_{\mathcal{H}}\le R\}\), the nonlinearity satisfies
\[
\|\mathcal{N}\Psi - \mathcal{N}\Phi\|
\le
L_R \|\Psi - \Phi\|
\]
for some finite constant \(L_R\) depending on \(R\), on \(n\), and on \(\max_k|\eta_k|\).
\end{lemma}

\begin{proof}
Work coordinatewise in a basis where \(f_n\) is componentwise. For scalar \(x,y\in\mathbb{C}\),
\(
|x^n - y^n|\le n \max(|x|,|y|)^{n-1}|x-y|.
\)
On the ball of radius \(R\) in \(\mathcal{H}\), each component is bounded by a constant \(C_R\), so \(f_n\) has Lipschitz constant at most \(n C_R^{n-1}\). Multiplying by \(|\eta_k|\le \eta_{\max}\) and summing over coordinates gives a global bound \(L_R\).
\end{proof}

For stability and well-posedness, it suffices that \(\mathcal{N}\) be locally Lipschitz; no global bound on \(\|\mathcal{N}\|\) is needed.

\subsection{Summary bound for the linear part}

Define the linear part
\[
\mathcal{L}_{\mathrm{lin}} \coloneqq \Xi + \Lambda_m + \mathcal{T} + \mathcal{D}.
\]
Combining Lemmas \ref{lem:Xi-norm}, \ref{lem:Lm-norm}, \ref{lem:D-norm}, and \ref{lem:T-norm} we have:

\begin{proposition}[Operator norm bound]
\label{prop:Llin-bound}
Suppose \(|\alpha_k|\le \alpha_{\max}\), \(\|B_k\|\le b_{\max}\) for all \(k\), \(|\lambda_k|\le \lambda_{\max}\) for all \(k\), \(\|W\|\le w\), and \(\|C\|\le c\). Then
\[
\|\mathcal{L}_{\mathrm{lin}}\|
\le
\alpha_{\max} + b_{\max} + w c + \lambda_{\max}\|L\|_2.
\]
\end{proposition}

\begin{proof}
By the triangle inequality for operator norms,
\[
\|\mathcal{L}_{\mathrm{lin}}\|
\le
\|\Xi\| + \|\Lambda_m\| + \|\mathcal{T}\| + \|\mathcal{D}\|,
\]
and each term is bounded by the corresponding lemma.
\end{proof}

\section{Existence and Stability of the Flow}

\subsection{Finite-dimensional ODE well-posedness}

Consider the deterministic evolution equation on \(\mathcal{H}\):
\begin{equation}
\frac{d\Psi}{dt} = \mathcal{L}_{\mathrm{lin}} \Psi + \mathcal{N}\Psi.
\label{eq:ODE-abstract}
\end{equation}
Since \(\mathcal{H}\) is finite-dimensional and \(\mathcal{N}\) is locally Lipschitz by Lemma \ref{lem:N-local}, standard ODE theory yields:

\begin{theorem}[Local existence and uniqueness]
For any initial state \(\Psi_0\in\mathcal{H}\), there exists a nonzero time \(T>0\) and a unique solution \(\Psi:[0,T)\to \mathcal{H}\) of \eqref{eq:ODE-abstract} with \(\Psi(0)=\Psi_0\).
\end{theorem}

Global existence on \([0,\infty)\) follows if we can control growth at infinity. One sufficient condition is a one-sided linear growth bound.

\begin{assumption}[Linear growth]
There exist constants \(a\ge 0\), \(b\ge 0\) such that for all \(\Psi\in\mathcal{H}\),
\[
\langle \Psi, \mathcal{L}_{\mathrm{lin}}\Psi\rangle_{\mathcal{H}} \le a\|\Psi\|_{\mathcal{H}}^2,
\quad
\|\mathcal{N}\Psi\|_{\mathcal{H}}\le b(1+\|\Psi\|_{\mathcal{H}}).
\]
\end{assumption}

\begin{theorem}[Global existence under linear growth]
Under the linear growth assumption, the solution of \eqref{eq:ODE-abstract} exists uniquely for all \(t\ge 0\).
\end{theorem}

\begin{proof}[Sketch]
The bound on \(\langle \Psi,\mathcal{L}_{\mathrm{lin}}\Psi\rangle\) gives an energy inequality of the form
\[
\frac{d}{dt}\|\Psi\|^2 \le 2a \|\Psi\|^2 + 2\|\Psi\|\,\|\mathcal{N}\Psi\|
\le 2a\|\Psi\|^2 + 2b\|\Psi\| + 2b\|\Psi\|^2.
\]
This is a Riccati-type inequality in \(\|\Psi\|\); Grönwall’s lemma yields a finite upper bound on \(\|\Psi(t)\|\) on any finite interval, preventing blow-up and extending the local solution globally.
\end{proof}

\subsection{Discrete-time Euler scheme and contractivity}

The MultiplicityCell implements a forward Euler step:
\begin{equation}
\Psi_{n+1} = \Psi_n + \Delta t\,F(\Psi_n),
\qquad F \coloneqq \mathcal{L}_{\mathrm{lin}} + \mathcal{N},
\label{eq:Euler-step}
\end{equation}
with fixed step size \(\Delta t>0\).

On a bounded subset where \(\mathcal{N}\) is Lipschitz with constant \(L_R\), we can bound the Lipschitz constant of \(F\) as
\[
\|F(\Psi)-F(\Phi)\|
\le
\|\mathcal{L}_{\mathrm{lin}}(\Psi-\Phi)\| + \|\mathcal{N}\Psi - \mathcal{N}\Phi\|
\le
\Bigl(\|\mathcal{L}_{\mathrm{lin}}\| + L_R\Bigr)\,\|\Psi-\Phi\|.
\]

\begin{proposition}[Euler step non-expansiveness]
\label{prop:euler-Lipschitz}
Assume that on a ball of radius \(R\), \(F\) is Lipschitz with constant \(L_R\). Then the Euler map \(\Phi(\Psi) = \Psi + \Delta t\,F(\Psi)\) is Lipschitz on that ball with constant
\[
\mathrm{Lip}(\Phi) \le 1 + \Delta t\, L_R.
\]
In particular, if \(L_R\ge 0\), the map is nonexpansive only if \(\Delta t = 0\); it is a contraction if \(L_R<0\) and \(\Delta t\) is small enough that \(1 + \Delta t L_R<1\).
\end{proposition}

\begin{proof}
For any \(\Psi,\Phi\),
\[
\|\Phi(\Psi)-\Phi(\Phi)\|
=
\|\Psi-\Phi + \Delta t (F(\Psi)-F(\Phi))\|
\le
\|\Psi-\Phi\| + \Delta t\,\|F(\Psi)-F(\Phi)\|
\le
(1+\Delta t L_R)\|\Psi-\Phi\|.
\]
\end{proof}

This makes explicit that stability of the discrete simulator is controlled by the operator norm of the linear part and the local Lipschitz constant of the nonlinearity, together with the step size \(\Delta t\).

\section{Norm Bounds for the Concrete MultiplicityCell}

In the implementation, the linear part of the cell comprises:

\begin{itemize}[nosep]
\item Prime evolution: diagonal with parameters \(\alpha_k\).
\item Scalar network: a small MLP applied per prime (linearized as \(B_k\)).
\item Mixing matrix: \(C\in\mathbb{R}^{N\times N}\).
\item Coupling projection: \(W:\mathbb{R}^d\to \mathbb{R}^d\).
\item Laplacian: fixed \(L\) on \(\mathbb{R}^N\).
\end{itemize}

Under the hypotheses of Proposition \ref{prop:Llin-bound}, we obtain:

\begin{corollary}[Practical upper bound]
Let
\[
\alpha_{\max} = \max_k|\alpha_k|,
\quad
b_{\max} = \max_k\|B_k\|,
\quad
\lambda_{\max} = \max_k|\lambda_k|.
\]
Let \(w = \|W\|\), \(c = \|C\|_2\), and \(\rho_L = \|L\|_2\). Then for the linear part implemented in the MultiplicityCell,
\[
\|\mathcal{L}_{\mathrm{lin}}\|
\le
\alpha_{\max}
+ b_{\max}
+ w c
+ \lambda_{\max} \rho_L.
\]
Consequently, on any bounded region where the nonlinearity is Lipschitz with constant \(L_R\), the combined vector field \(F=\mathcal{L}_{\mathrm{lin}}+\mathcal{N}\) has Lipschitz constant at most
\[
L_R^{\mathrm{tot}} \le \alpha_{\max}
+ b_{\max}
+ w c
+ \lambda_{\max} \rho_L
+ L_R,
\]
and the Euler step with step size \(\Delta t\) has Lipschitz constant at most
\[
\mathrm{Lip}(\Phi) \le 1 + \Delta t\,L_R^{\mathrm{tot}}.
\]
\end{corollary}

These bounds provide explicit stability criteria in terms of the learnable parameters:
\begin{itemize}[nosep]
\item If the spectral norms of \(B_k,W,C\) and the magnitudes of \(\alpha_k,\lambda_k\) are regularized during training, one can ensure that \(L_R^{\mathrm{tot}}\) stays in a controlled range.
\item The step size \(\Delta t\) then trades off between approximation quality of the continuous flow and stability of the discrete simulator.
\end{itemize}

\section{Remarks on Extensions}

\subsection{Memory integral as auxiliary state}

The original UME includes a memory term \(\beta_k \int_0^t I_k(\tau)\,d\tau\). A mathematically clean way to include it in this operator framework is to augment the state by an auxiliary component \(h_k(t)\) satisfying
\[
\frac{d h_k}{dt} = I_k(\Psi(t),t),
\qquad
(\mathcal{M}\Psi)_k = \beta_k h_k.
\]
The extended state space is \(\mathcal{H}\times\mathcal{H}\); the above norm and Lipschitz arguments extend verbatim to this larger finite-dimensional ODE system.

\subsection{Geometric terms}

If a Berry-connection or Fubini--Study based term is inserted directly into the generator rather than as a penalty, one obtains additional contributions to \(\mathcal{L}_{\mathrm{mult}}\) of the schematic form
\[
\mathcal{G}[\Psi] = G(\Psi)\Psi
\]
for some operator-valued functional \(G\). So long as \(G\) is locally bounded and locally Lipschitz in \(\Psi\), the existence, uniqueness, and Euler-step Lipschitz results remain valid. The nontrivial analysis would then be to bound \(\|G(\Psi)\|\) in terms of \(\|\Psi\|\); this is left for future work.

\section{Conclusion of the Appendix}

The purpose of this appendix has been to make explicit:

\begin{itemize}[nosep]
\item The operator-norm bounds and Lipschitz constants for each contribution to the generator \(\mathcal{L}_{\mathrm{mult}}\) in the finite-dimensional multiplicity setting.
\item The resulting well-posedness of the continuous-time flow and the stability properties of the discrete Euler simulator used in the MultiplicityCell.
\end{itemize}

These results do not by themselves validate the modeling assumptions of the UME, but they provide a mathematically controlled foundation on which to tune parameters and step size, and to reason about the qualitative stability of the learned dynamics.

\nocite{*}
\bibliographystyle{unsrt}  
\bibliography{references}
\end{document}